\subsection{Polynomial mappings}

DEFINITION 4. --- Let M and N be A-modules and assume that M is free. Let \(\operatorname{Map}(M, N)\) be the A-module of all mappings of M into N. The submodule \(\sum_{q > 0} \operatorname{Pol}_{A}^{q}(M, N)\) of \(\operatorname{Map}(M, N)\) is denoted by \(\operatorname{Pol}_{A}(M, N)\) or simply \(\operatorname{Pol}(M, N)\) ;

its elements are called polynomial mappings of M into N.

Let \((e_{i})_{i\in I}\) be a basis of M and suppose that I is finite; by Prop. 13 (IV, p.~54) a mapping f of M into N is polynomial if and only if there exists a polynomial F in the indeterminates \(X_{i}\) with coefficients in N such that

\[
f \left(\sum_ {i \in I} x _ {i} e _ {i}\right) = F (\mathbf {x})
\]

for every family \(\mathbf{x} = (x_i)_{i\in I}\) in \(\mathbf{A}^{(I)}\) . This property is independent of the basis chosen for \(\mathbf{M}\) and it justifies the terminology « polynomial mapping ».

PROPOSITION 17. --- Let M be a free A-module and B an associative, commutative and unital A-algebra. Then \(\operatorname{Pol}_{\mathbf{A}}(\mathbf{M}, \mathbf{B})\) is a sub-B-algebra of the algebra \(\operatorname{Map}(\mathbf{M}, \mathbf{B})\) .

This follows from Def. 4 and Prop. 13, (iv) (IV, p.~54).

PROPOSITION 18. --- Let M, N, P be A-modules, and assume that M and N are free. If \(f \in \operatorname{Pol}(M, N)\) , \(g \in \operatorname{Pol}(N, P)\) , then \(g \circ f \in \operatorname{Pol}(M, P)\) .

We can reduce at once to the case where there exists an integer q such that \(g \in \text{Pol}^{q}(N, P)\) ; then there exists a q-linear mapping h of \(N^{q}\) into P such that \(g(y) = h(y, y, \ldots, y)\) for all \(y \in N\) . Writing f as a sum of homogeneous polynomial mappings we are thus reduced to proving that the mapping

\[
x \mapsto h (f _ {1} (x), f _ {2} (x), \dots , f _ {q} (x))
\]

of M into P, where \(f_{i} \in \operatorname{Pol}^{q_{i}}(\mathbf{M}, \mathbf{N})\) , is polynomial. For \(i = 1, \ldots, q\) there exists a \(q_{i}\) -linear mapping \(l_{i}\) of \(M^{q_{i}}\) into N such that \(f_{i}(x) = l_{i}(x, x, \ldots, x)\) for all \(x \in M\) . Hence

\[
h \left(f _ {1} (x), f _ {2} (x), \dots , f _ {q} (x)\right) = h \left(l _ {1} (x, \dots , x), \dots , l _ {q} (x, \dots , x)\right),
\]

from which our assertion follows.

Lemma 2. --- Let N be an A-module, n an integer ≥0 and

\[
f = m _ {0} + m _ {1} \mathrm{X} + \dots + m _ {n} \mathrm{X} ^ {n} \in \mathrm{N} [ \mathrm{X} ].
\]

Suppose that there exist \(\alpha_{0}, \alpha_{1}, \ldots, \alpha_{n} \in A\) such that \(f(\alpha_{0}) = \cdots = f(\alpha_{n}) = 0\) , and such that for \(i \neq j\) the homothety of ratio \(\alpha_{i} - \alpha_{j}\) in N is injective, then f = 0.

(This lemma generalizes the Cor. of IV, p.~16.)

The lemma clearly holds for \(n = 0\) ; we shall prove it by induction on \(n\) . We have

\[
f (\mathbf {X}) = f (\mathbf {X}) - f \left(\alpha_ {0}\right) = \sum_ {i = 1} ^ {n} m _ {i} \left(\mathbf {X} ^ {i} - \alpha_ {0} ^ {i}\right) = \left(\mathbf {X} - \alpha_ {0}\right) g (\mathbf {X})
\]

where g is an element of N{[}X{]} of the form \(m_{0}^{\prime} + m_{1}^{\prime} X + \cdots + m_{n-1}^{\prime} X^{n-1}\) . The hypotheses of the lemma imply that \(g(\alpha_{1}) = \cdots = g(\alpha_{n}) = 0\) , hence g = 0 by the induction hypothesis, and so f = 0.

PROPOSITION 19. --- Let M be a free A-module, N an A-module, G an infinite additive subgroup of A, and suppose that the homotheties of N defined by the nonzero elements of G are injective. Then \(\operatorname{Pol}(M,N)\) is the direct sum of the \(\operatorname{Pol}^{q}(M,N)\) .

Let \(f_{0}, f_{1}, \ldots, f_{n}\) be such that \(f_{i} \in \operatorname{Pol}^{i}(\mathbf{M}, \mathbf{N})\) and suppose that we have the relation \(f_{0} + \cdots + f_{n} = 0\) . Let \(x \in \mathbf{M}\) , then for all \(\lambda \in G\) we have

\[
0 = \sum_ {i = 0} ^ {n} f _ {i} (\lambda x) = \sum_ {i = 0} ^ {n} \lambda^ {i} f _ {i} (x).
\]

By Lemma 2, applied to the polynomial \(\sum_{i=0}^{n} f_i(x) X^i\) we have

\[
f _ {0} (x) = \dots = f _ {n} (x) = 0.
\]

COROLLARY. --- Assume that A is an infinite integral domain; let M be a free A-module and N a torsion-free A-module.

\begin{enumerate}
\def\labelenumi{(\roman{enumi})}
\tightlist
\item
  We have \(\operatorname{Pol}(\mathbf{M},\mathbf{N}) = \bigoplus_{q\geqslant 0}\operatorname{Pol}^q (\mathbf{M},\mathbf{N})\) and each \(\operatorname{Pol}^q (\mathbf{M},\mathbf{N})\) may be
\end{enumerate}

identified canonically with \(\operatorname{Hom}(\mathbf{TS}^q(\mathbf{M}),\mathbf{N})\) .

\begin{enumerate}
\def\labelenumi{(\roman{enumi})}
\setcounter{enumi}{1}
\tightlist
\item
  Let \(f \in \operatorname{Pol}(M, N)\) and \((e_i)_{i \in I}\) a basis of \(M\) . There exists one and only one family \((u_\nu)_{\nu \in \mathbb{N}^{(I)}}\) of elements of \(N\) such that \(f\left(\sum_{i \in I} \lambda_i e_i\right) = \sum_{\nu \in \mathbb{N}^{(I)}} \lambda^\nu u_\nu\) for all
\end{enumerate}

\[
\left(\lambda_ {i}\right) \in \mathbf {A} ^ {(I)}
\]

The assertion (i) follows from Prop. 16 and 19, and (ii) follows from (i) and the Cor. of Prop. 16.
% semantic-fragments: legacy-input-seam
\relax{}
